\chapter{Ameaças à validade e implicações práticas}
\label{cap:ameacas-validade-implicacoes-praticas}

Este capítulo apresenta as principais ameaças à validade do trabalho e discute suas implicações práticas. As limitações da avaliação, das métricas utilizadas e dos mecanismos de guardrail são tratadas dentro dessas ameaças, evitando a repetição de pontos já discutidos no capítulo de resultados. Dessa forma, o capítulo delimita o alcance dos achados e indica os cuidados necessários para interpretar a solução dentro de seu escopo funcional, linguístico e exploratório.

% \section{Limitações}
% \label{sec:limitacoes-escopo-solucao}

% Como delimitado na proposta, a solução foi concebida como um mecanismo de adaptação comunicacional de textos técnicos de saúde previamente definidos ou validados. Sua atuação se concentra na forma de apresentação da informação, e não na definição do conteúdo clínico. Por isso, os resultados da avaliação permitem observar o comportamento da simplificação textual, da adaptação por perfil e dos guardrails, mas não permitem inferir eficácia clínica, segurança de uso em contexto assistencial real ou impacto sobre a adesão ao tratamento.

% Também não foi avaliado se os textos gerados melhoram a execução de cuidados, reduzem erros em tratamentos ou produzem efeitos concretos sobre desfechos de saúde. Esses aspectos exigiriam estudos específicos, com participação de usuários reais e validação por profissionais da saúde. Assim, os resultados devem ser compreendidos dentro de um escopo funcional e linguístico, não como evidência de validação clínica dos textos produzidos.

% \subsection{Limitações da avaliação}
% \label{sec:limitacoes-avaliacao}

% A avaliação realizada utilizou dois materiais de saúde, três perfis simulados de pacientes e 18 execuções principais. Essa configuração foi adequada para observar o comportamento inicial da solução, comparar diferentes níveis de adaptação comunicacional e analisar a atuação dos guardrails em um cenário controlado. No entanto, ela não permite generalizar os resultados para todos os tipos de documentos de saúde, perfis de pacientes ou contextos de uso em aplicações mHealth.

% Os dois materiais utilizados apresentam orientações práticas e linguagem técnica, mas não representam toda a diversidade de conteúdos encontrados em saúde digital. Materiais envolvendo medicamentos, dosagens, horários, contraindicações detalhadas, sinais de alerta, protocolos clínicos complexos ou orientações pós-operatórias poderiam apresentar desafios diferentes. Nesses casos, pequenas alterações textuais poderiam ter impacto mais sensível sobre a interpretação das orientações.

% Da mesma forma, os três perfis simulados permitiram observar variações de escolaridade e familiaridade com linguagem técnica, mas não cobrem a diversidade real de usuários. Pacientes podem variar quanto à idade, letramento em saúde, experiência prévia com a condição, limitações cognitivas, preferências de linguagem e contexto social. Além disso, como não houve participação de pacientes reais nem testes de compreensão com usuários finais, não é possível afirmar que os textos simplificados seriam efetivamente compreendidos com mais facilidade em situações reais de uso.

% \subsection{Limitações das métricas de legibilidade}
% \label{sec:limitacoes-metricas-legibilidade}

% O Flesch-PT foi utilizado como métrica automática de legibilidade por permitir uma comparação objetiva entre os textos originais e as versões simplificadas. No contexto da avaliação, ele foi útil para indicar mudanças gerais do texto, como possíveis reduções de complexidade associadas à estrutura das frases e às características das palavras utilizadas. No entanto, essa métrica representa apenas uma dimensão limitada da legibilidade.

% O índice não mede compreensão real do usuário, clareza percebida, adequação comunicacional completa ou precisão médica. Um texto pode apresentar Flesch-PT mais alto e ainda assim conter problemas de organização, ambiguidade, perda de informação relevante ou inadequação ao perfil do leitor. Da mesma forma, um texto pode manter Flesch-PT baixo por preservar termos técnicos necessários, sem que isso signifique necessariamente uma pior adaptação.

% O caso do perfil Ana no material sobre diabetes ilustra essa limitação. Nesse cenário, a variação média do Flesch-PT foi negativa, mas Ana representa uma usuária com ensino superior e familiaridade com a área da saúde. Portanto, a manutenção de terminologia mais técnica pode ser coerente com o perfil comunicacional definido. Esse resultado mostra que a qualidade da adaptação não pode ser avaliada apenas pelo aumento do índice, pois a simplificação adequada depende também do público-alvo e da necessidade de preservar determinados termos.

% Dessa forma, o Flesch-PT deve ser interpretado como um indicador auxiliar, e não como uma medida completa de qualidade textual. Sua utilização contribui para a análise exploratória da solução, mas precisa ser complementada por avaliação qualitativa, revisão e, em trabalhos futuros, testes de compreensão com usuários finais.

% \subsection{Limitações dos guardrails}
% \label{sec:limitacoes-guardrails}

% Os guardrails implementados contribuíram para identificar alterações potencialmente problemáticas nas versões geradas, como troca de substância, alteração de instrução procedimental, alteração de sentido semântico, contradição numérica e possível perda de negação. Ainda assim, eles não garantem preservação absoluta do conteúdo original. A existência de uma camada de verificação reduz riscos, mas não elimina a possibilidade de falhas.

% O guardrail baseado em LLM depende da capacidade do próprio modelo de comparar o texto original e a versão simplificada. Esse tipo de verificação pode ser sensível à formulação do prompt, ao modelo utilizado, à configuração de geração e à complexidade do trecho analisado. Em alguns casos, o modelo pode deixar de identificar uma alteração relevante ou interpretar como problemática uma reformulação aceitável.

% O checker determinístico, por sua vez, possui limitações diferentes. Por ser baseado em regras e expressões regulares, ele depende dos padrões previamente definidos na implementação. Esse tipo de verificação pode ser rígido demais, gerando falsos positivos quando uma informação é preservada por meio de outra construção linguística, mas também pode deixar de detectar casos que não se encaixam nos padrões previstos. Isso é especialmente relevante em linguagem natural, na qual negações, instruções e equivalências semânticas podem aparecer de formas variadas.

% Apesar dessas limitações, a postura conservadora adotada é adequada ao domínio de saúde. Em textos que envolvem orientações de cuidado, é preferível rejeitar uma saída e gerar uma nova tentativa do que apresentar ao usuário uma versão com possível alteração de sentido, perda de negação ou mudança em uma instrução relevante. Portanto, os guardrails devem ser compreendidos como mecanismos de mitigação de risco, e não como garantia definitiva de correção clínica.

\section{Ameaças à validade}

As ameaças à validade deste trabalho estão relacionadas às escolhas de implementação, ao escopo controlado da avaliação e às formas utilizadas para medir os conceitos analisados. A seguir, são sintetizadas as principais ameaças identificadas.

\begin{itemize}
    \item \textbf{Validade interna:} A validade interna pode ter sido influenciada pelo comportamento não determinístico da LLM, pelas configurações do modelo, pela formulação dos prompts e pelos critérios adotados nos guardrails. Mesmo com instruções definidas, pequenas variações entre execuções podem afetar o grau de simplificação, a escolha lexical e a organização do texto. Além disso, as características dos dois documentos avaliados, como tema, vocabulário e presença de instruções restritivas, também podem ter influenciado os valores de Flesch-PT e a ocorrência de rejeições pelos guardrails.
    
    \item \textbf{Validade externa:} Os resultados não podem ser generalizados diretamente para outros materiais de saúde, perfis de pacientes, modelos de linguagem, idiomas ou contextos clínicos. A avaliação foi realizada em português, com dois documentos específicos, três perfis simulados e uma implementação acadêmica. Conteúdos com maior criticidade clínica, como materiais envolvendo medicamentos, dosagens, contraindicações ou protocolos complexos, poderiam exigir mecanismos de verificação mais robustos e revisão profissional mais rigorosa. Da mesma forma, a adoção da solução em uma plataforma mHealth real, como a Takere, exigiria novos testes, integração operacional e validação por profissionais responsáveis pelos conteúdos.
    
    \item \textbf{Validade de construto:} A validade de construto está relacionada à forma como conceitos como legibilidade, simplificação, preservação das informações essenciais e adequação ao perfil foram avaliados. A legibilidade foi medida pelo Flesch-PT, que considera aspectos formais do texto, mas não mede compreensão real, clareza percebida ou esforço cognitivo. A preservação das informações essenciais foi analisada por leitura manual preliminar e verificação automática, o que não equivale a validação clínica independente. Além disso, os perfis utilizados foram simulados e não representam toda a diversidade de usuários reais, que podem variar quanto a letramento em saúde, experiência prévia, preferências de linguagem e contexto social.
\end{itemize}

\section{Implicações práticas}
\label{sec:sintese-limitacoes}

As ameaças apresentadas não invalidam os resultados obtidos, mas delimitam seu alcance. A avaliação realizada fornece evidências iniciais de que a solução é capaz de simplificar textos técnicos de saúde na maior parte dos casos avaliados, adaptar a linguagem conforme perfis comunicacionais simulados e utilizar guardrails para bloquear alterações potencialmente problemáticas. No entanto, esses achados devem ser compreendidos em caráter funcional, linguístico e exploratório.

O trabalho não demonstrou validação clínica, compreensão efetiva por pacientes reais ou impacto em adesão ao tratamento. Além disso, a avaliação foi limitada a dois materiais, três perfis simulados, uma métrica automática de legibilidade e uma análise manual preliminar. Como discutido anteriormente, os guardrails contribuíram para reduzir riscos, mas não garantem preservação absoluta do conteúdo original.

Dessa forma, o uso da abordagem em contexto real exigiria validação mais ampla, participação de profissionais da saúde, avaliação com usuários finais, ampliação dos materiais testados e aprimoramento contínuo dos mecanismos de verificação. Ainda assim, dentro do escopo proposto, os resultados apoiam a viabilidade da solução como apoio à adaptação comunicacional de textos de saúde previamente definidos, desde que seu uso seja mantido dentro de limites claros e, em contextos reais, seja acompanhado por processos adequados de revisão e validação.